\section{Gruppearbejde}






% jonathan
Jeg tog en mindre rolle i projektets opstart og dette påvirkede mit engagement i de tidlige faser af projektet, og jeg oplevede også perioder, hvor jeg ikke deltog i møder eller informerede gruppen om fravær på forhånd. Set i bakspejlet kunne jeg have bidraget mere aktivt. Jeg tror, at min manglende deltagelse i starten kan have været en form for social loafing, hvor jeg ubevidst trak mig tilbage fra ansvaret, fordi jeg følte, at andre i gruppen ville tage sig af opgaverne, eller allerede havde taget sig af dem. For at bryde ud af denne dynamik forsøgte jeg at bidrage på områder, hvor jeg oplevede, at der manglede arbejdskraft, blandt andet gennem kildearbejde og rapportskrivning. Da jeg som person har en tendens til at være tilbageholdene i sociale sammenhænge, var dette en måde jeg stadig kunne bidrage på.
Jeg oplevede at gruppen håndterede situationen professionelt og inkluderende, hvilket gjorde det lettere for mig at bidrage mere aktivt senere i projektforløbet. Jeg har lært, at det er vigtigt at kommunikere åbent om ens deltagelse og forpligtelser i et gruppeprojekt for at sikre, at alle medlemmer føler sig ansvarlige og engagerede. I fremtidige projekter vil jeg stræbe efter at være mere proaktiv i min deltagelse og kommunikere klart med mine gruppemedlemmer om mine bidrag og eventuelle udfordringer, jeg måtte have.